% BEGIN REV09_GERMEN_DODECAFASICO
\section{Dodecaphasic germ and recovery of unimodal branches}
\label{corpus9:iv:germen}

The oriented wheel \(U_{12}^{\mathrm{sign}}\), received from the
TPK return, determines, in the uniform sector \(\eta=0\),
the weights \(w_m=1/12\). Fix the phase functional and its ordered
real representatives
\begin{equation}
 \phi_m=\pi_{\mathrm{HMT}}\frac{3m-1}{18},
 \qquad m=1,\ldots,12.
 \label{corpus9:iv:germen-fases}
\end{equation}
The real lift of the phases is part of the reader:
the following moments are evaluated on these representatives.
The construction proceeds from the wheel, its orientation, weights
and functional; a modification of any of these data
modifies the resulting collection.

\subsection{Normalization and analytic profile}
Let \(a_m=(3m-1)/18\) and
\(D_0=\frac1{12}\sum_{m=1}^{12}a_m^2\). The integer sums
\[
 \sum_{m=1}^{12}(3m-1)^2=5394=6\cdot899,\qquad
 \sum_{m=1}^{12}(3m-1)^4=4294614
\]
give \(D_0=899/648\). The quadratic normalization condition
\(s_0^2\pi_{\mathrm{HMT}}^2D_0=2\) fixes
\[
 s_0=\frac{36}{\pi_{\mathrm{HMT}}\sqrt{899}},
 \qquad k_m=s_0\phi_m=\frac{2(3m-1)}{\sqrt{899}}.
\]
The angular scale cancels in this normalization before
iteration of the profile
\begin{equation}
 u_0(x)=1-\frac1{12}\sum_{m=1}^{12}\cos(k_mx),
 \qquad f_\lambda(x)=1-\lambda u_0(x).
 \label{corpus9:iv:germen-perfil}
\end{equation}

\begin{proposition}[Exact representation and criticality]
\label{corpus9:iv:germen-analitico}
The profile \(u_0\) is an even entire function and satisfies
\[
 u_0(x)=x^2-\frac{715769}{2424603}x^4+O(x^6).
\]
With \(y=x/\sqrt{899}\), its closed expression is
\begin{equation}
 u_0(x)
  =1-\frac{\cos(37y)\sin(36y)}{12\sin(3y)}
  =1-\frac{\sin(73y)-\sin y}{24\sin(3y)},
 \label{corpus9:iv:germen-cerrado}
\end{equation}
where the quotients extend through their removable singularities.
For \(0<\lambda\leq2/u_0(1)\), \(f_\lambda\) is an even,
real-analytic self-map of \([-1,1]\), with a unique
quadratic critical point at zero and negative Schwarzian
away from that point.
\end{proposition}
\begin{proof}
The finite cosine sum is entire and even. Its Taylor
expansion uses
\[
 \frac1{12}\sum_m k_m^2=2,\qquad
 \frac1{24\cdot12}\sum_m k_m^4=\frac{715769}{2424603},
\]
obtained from the preceding integer sums. For the closed
form, the progression of arguments is \((6m-2)y\). Summing
the geometric series of exponentials and taking the real part gives
\[
 \sum_{m=1}^{12}\cos((6m-2)y)
   =\frac{\sin(36y)}{\sin(3y)}\cos(37y).
\]
The identity \(2\cos(37y)\sin(36y)=\sin(73y)-\sin y\)
proves \eqref{corpus9:iv:germen-cerrado}. The finite sum defines
its extension at every zero of the denominator.

For \(0<x\leq1\), all arguments satisfy
\(0<k_mx\leq70/\sqrt{899}<\pi_{\mathrm{HMT}}\). Hence
\[
 u_0'(x)=\frac1{12}\sum_m k_m\sin(k_mx)>0,\qquad
 u_0'''(x)=-\frac1{12}\sum_m k_m^3\sin(k_mx)<0.
\]
Parity gives the opposite signs on the negative interval
and proves uniqueness of the critical point. Moreover,
\(f_\lambda''(0)=-2\lambda\ne0\). Since
\(0\leq u_0(x)\leq u_0(1)\) on \([-1,1]\), the stated
interval for \(\lambda\) gives
\(-1\leq f_\lambda(x)\leq1\). Finally,
\[
 S(f_\lambda)
  =\frac{u_0'''}{u_0'}
   -\frac32\left(\frac{u_0''}{u_0'}\right)^2<0
 \qquad(0<|x|\leq1).
\]
The first quotient is negative on both sides of the origin,
and the remaining term is nonpositive.
\end{proof}

\subsection{Branchwise inversion and itinerary memory}
\begin{proposition}[Recovery of the folding]
\label{corpus9:iv:germen-inversa}
Let \(\lambda>0\) and
\(v=(u_0|_{[0,1]})^{-1}\). For
\(y\in[1-\lambda u_0(1),1]\), the preimage on a branch is
\begin{equation}
 x=\sigma\,v\!\left(\frac{1-y}{\lambda}\right),
 \qquad \sigma=\operatorname{sgn}x.
 \label{corpus9:iv:germen-rama}
\end{equation}
At \(y=1\), the preimage is zero and \(\sigma=0\) is fixed.
\end{proposition}
\begin{proof}
Continuity and strict monotonicity of \(u_0\) on
\([0,1]\) provide the inverse \(v\).
The equation \(y=1-\lambda u_0(x)\), together with parity,
determines \(|x|=v((1-y)/\lambda)\); the sign determines
orientation. At \(y=1\), uniqueness of the minimum of \(u_0\)
gives \(x=0\). Away from the critical point, the inverse is
differentiable. Its derivative becomes singular at the
origin, where \(u_0(x)=x^2+O(x^4)\).
\end{proof}

For a finite distribution on pairs \(\pm x_i\), with distinct
\(x_i>0\), let \(p_i^\pm\) be their probabilities,
\(p_i=p_i^++p_i^-\), \(Y=f_\lambda(X)\) and
\(\Sigma=\operatorname{sgn}X\). A mass at zero is also
allowed; its branch-uncertainty contribution is zero. With
natural logarithms,
\begin{align}
 H(X\mid Y,\Sigma)&=0,\\
 H(X\mid Y)&=H(\Sigma\mid Y)
   =\sum_{i:p_i>0}p_i\,h(p_i^+/p_i),\\
 h(p)&=-p\log p-(1-p)\log(1-p).
 \label{corpus9:iv:germen-entropia}
\end{align}
The proof groups the distribution by the fibres of
\(f_\lambda\). Monotonicity distinguishes the moduli \(x_i\),
and \eqref{corpus9:iv:germen-rama} distinguishes the two signs
within each fibre. With equal probabilities in each pair
and no mass at zero, the uncertainty is \(\log2\).

If \(0<\lambda\leq2/u_0(1)\), the second iteration is
defined on the same interval and satisfies
\[
 H(X\mid f_\lambda^2(X))
 =H(X\mid f_\lambda(X))
  +H(f_\lambda(X)\mid f_\lambda^2(X)).
\]
Indeed, both intermediate maps are deterministic,
and the chain rule for the entropy of
\((X,f_\lambda(X))\), conditioned on \(f_\lambda^2(X)\),
gives the identity. Successive recovery uses the signs
at each step. For \(\lambda=0\), the map is constant
and additionally requires retention of the input modulus.

\subsection{Scope of the rescaled return}
For \(a=-f(1)\ne0\), let \(h_a(x)=-ax\). The operator
\[
 \mathcal R(f)=h_a^{-1}\circ f^2\circ h_a
\]
is a rescaled return on \([-1,1]\) when the interval
\(J=h_a([-1,1])\) belongs to the domain of \(f^2\) and
satisfies \(f^2(J)\subseteq J\). The scaling \(h_a\) is
bijective; the two-step return retains its interpretation
with the itinerary and the map \(f\). Reconstruction of an
iterative root from \(\mathcal R(f)\) and \(a\) is an
additional condition. Proposition
\ref{corpus9:iv:germen-analitico} establishes the germ and
its critical class. Existence and uniqueness of a fixed point
of the renormalizer, its stable and unstable spectral
decomposition, and transversality of this collection are the
additional hypotheses required for passage to a universal
renormalization limit.
% END REV09_GERMEN_DODECAFASICO
